\chapter{Just a bunch of syllables\ldots}

\section{Groove solkattu}

Solkattu (Tamil) translates literally to 'bunch of syllables'. These syllables are
used to count-out or recite the various patterns. This aids practice.

Let us begin with the various terms used to describe walks/grooves. There is a space
after every atom. A hyphen (-) denotes a silent syllable (sometimes called swallowed
syllable or \emph{kaarvai}): the sol or syllable exists, but is not recited.

\begin{center}
\begin{adjustbox}{max width=\textwidth}
\begin{tabular}{l c}
\toprule
Sol & atoms\\
\midrule
Tha ki ta                                  & 3\\
Tha ka dhi mi                              & 4\\
Tha ka tha ki ta                           & 5\\
Dhin - - tha ki ta                         & 6\\
Tha ka dhi mi tha ki ta                    & 7\\
Tha ka dhi mi tha ka ju nu                 & 8\\
Tha ka dhi mi, tha ka tha ki ta (4 + 5)    & 9\\
\bottomrule
\end{tabular}
\end{adjustbox}
\end{center}

There is another group of sol\,s used as conclusive fillers. Note that although the
length of some of these sol\,s is the same as the ones above, the purpose is
different. Conclusive fillers are used at the end of a pattern. The different types
of patterns are dealt with in the next chapter. A comma indicates conjunction and
nothing else. There is no pause.

\begin{center}
\begin{adjustbox}{max width=\textwidth}
\begin{tabular}{l c}
\toprule
Sol & atoms\\
\midrule
Tha dhi ki ta thom                  & 5\\
Tha dhi tha kita thom               & 5\\
Kita thaka thari kita thom          & 5\\
Tha dhee - ki ta thom               & 6\\
Tha - dhee - ki ta thom             & 7\\
Tha dhin - gin - na - thom          & 8\\
Tha thom - tha dhi ki ta thom       & 8\\
Tha ka, tha dhee - ki ta thom       & 8\\
Tha - dhee - ki - ta - thom         & 9\\
Tha ka dhi na, tha dhi ki ta thom   & 9\\
Tha ki ta thom -, tha dhi ki ta thom & 10\\
Tha ka tha ka, tha dhee - ki ta thom & 10\\
\bottomrule
\end{tabular}
\end{adjustbox}
\end{center}

\section{More grooves and fillers}

\subsection{Grooves}

Below is a list of grooves. Equally spaced. 8 atoms each:
\begin{enumerate}
  \item tha na, tha dhi nu, tha dhi nu
  \item tha na, dhi nu, dhi nu, na tha
  \item tha dhi nu, tha na, tha dhi nu
  \item Dhi nu nam thari kita tha dhi nu
\end{enumerate}

Below is a list of grooves. Equally spaced. 6 atoms each. The last one is 12 atoms
long:
\begin{enumerate}
  \item than-gu tha tha dhi na (thangu is two atoms long).
  \item tha kita thaka dhin dhin na
  \item dhin dhin na dhin dhin na
  \item dhi nu nam thari kita nam dhi nu dhi nu na tha (12 atoms)
\end{enumerate}

This one is ten atoms long: dhi nu nam thari kita nam dhi nu na tha. Associated
filler: nam thari kita nam - nam - nam - tha. Now, can you find the associated filler
for No.~4 from the set above? Hint: look at the similarity of the 12 matra version to
this one. Each nam stands for 2 matras --- note the hyphen after the nam. This one is
ten atoms too: Tha na tha dhi nu tha, tha na dhi nu. Note the similarity to No.~1 from
the 8-atom set above.

\subsection{Fillers}

\begin{enumerate}
  \item tham kita thaka dhina (4)
  \item thaka thari kita thaka (4)
  \item tha langu thaka dhina (4)
  \item nam thom kita thaka (4)
  \item tha dhith thaka dhina (4)
  \item tha tha ki ta tha ka (4)
  \item nam thari kita dhin thari kira tha langu (8)
  \item tha kita thaka dhina naka dhith tha langu (8)
  \item tham kita thaka thari kita thaka (6)
  \item thari kita thaka thari kita thaka (6)
  \item thari jina kita jinu thaka thari kita thaka (8)
  \item No.~4 + No.~10 = tha tha kita thaka, thari kita, thaka thari kita thaka (10)
\end{enumerate}

Note that some of the above sols are actually twice as long as mentioned. For example,
the last one can be thought of as being Tha - tha - ki ta tha ka, tha ri ki ta, tha ka
tha ri ki ta tha ka. That would be 20 atoms long. However, playing it at twice the
speed (double time) is more common.
